\chapter{Thoracoscopy}
\index{Thoracoscopy}

\medskip
\noindent\textit{Educational reference; always seek professional advice for individual care.}

\section*{Definition and Overview}
Thoracoscopy is a procedure that allows a doctor to look inside the pleural space---the thin cavity between the lung and the chest wall---using a rigid or flexible telescope inserted through small incisions between the ribs. Medical thoracoscopy (pleuroscopy) is often performed by respiratory physicians under local anaesthesia and sedation to diagnose pleural effusion and to perform talc pleurodesis. Surgical thoracoscopy, better known as video-assisted thoracoscopic surgery (VATS), is performed by thoracic surgeons under general anaesthesia with one-lung ventilation and can include lung biopsy, wedge resection, lobectomy, decortication, and treatment of pneumothorax. Compared with open thoracotomy, thoracoscopic approaches usually mean less pain, shorter hospital stay, and faster recovery for suitable patients. This chapter explains why thoracoscopy is offered, how to prepare, what happens during and after the procedure, and which warning signs need urgent attention.

\section*{Epidemiology}
Pleural disease is common in respiratory practice. Malignant pleural effusion from lung cancer, breast cancer, and mesothelioma; parapneumonic effusion and empyema; tuberculous pleuritis; and unexplained exudates frequently require tissue diagnosis. historical asbestos exposure keeps mesothelioma and related pleural disease relevant, especially in older men with occupational exposure. VATS has become standard for many spontaneous pneumothorax operations, early-stage lung cancer resections in fit patients, and some mediastinal procedures. Access depends on thoracic surgical and interventional pulmonology services, concentrated in larger hospitals, with referral pathways from regional centres.

\section*{Aetiology and Pathophysiology}
Thoracoscopy is indicated when imaging and fluid analysis leave important questions unanswered, or when a therapeutic intervention in the pleural space is required.

\begin{itemize}
    \item \textbf{Pleural effusion mechanisms:} imbalance of production and absorption from malignancy, infection, heart failure, inflammation, or trauma
    \item \textbf{Need for parietal pleural biopsy:} fluid cytology alone misses many malignancies and granulomatous diseases; direct biopsy raises diagnostic yield
    \item \textbf{Pleurodesis principle:} talc or other agents inflame and fuse visceral and parietal pleura to prevent re-accumulation of fluid or air
    \item \textbf{Pneumothorax paths:} bleb rupture and air leak; VATS staples blebs and may add pleurodesis
    \item \textbf{Lung nodule and mass assessment:} VATS allows targeted resection with histological margins when appropriate
    \item \textbf{Empyema stages:} loculated infected collections may need VATS drainage and decortication when chest tube drainage fails
\end{itemize}

\section*{Clinical Features}
People considered for thoracoscopy often have:

\begin{itemize}
    \item \textbf{Breathlessness and chest heaviness:} from moderate to large pleural effusion
    \item \textbf{Pleural pain and cough:} especially with infection or pneumothorax
    \item \textbf{Recurrent effusion after tapping:} suggesting need for definitive diagnosis and pleurodesis
    \item \textbf{Undiagnosed exudative effusion:} after blood tests, CT, and thoracentesis
    \item \textbf{Suspected mesothelioma or metastatic disease:} needing histology for treatment planning
    \item \textbf{Spontaneous pneumothorax:} first large or recurrent episodes in selected patients
    \item \textbf{Persistent air leak or empyema:} not resolving with conservative drainage
\end{itemize}

\section*{Differential Diagnosis}
Thoracoscopy is a tool, not a diagnosis. Conditions evaluated or treated include:

\begin{itemize}
    \item \textbf{Malignant versus benign effusion:} cancer, mesothelioma, versus inflammatory or cardiac causes
    \item \textbf{Tuberculous versus non-tuberculous pleural disease:} granulomas and culture guidance
    \item \textbf{Parapneumonic effusion versus established empyema:} determines urgency of drainage
    \item \textbf{Primary spontaneous versus secondary pneumothorax:} COPD and underlying lung disease change risk
    \item \textbf{Alternatives to thoracoscopy:} image-guided percutaneous biopsy, bronchoscopy, medical thoracentesis alone, or open surgery when ports cannot achieve the goal
\end{itemize}

\section*{When to Seek Medical Care}
Discuss new or worsening breathlessness, chest pain, or known effusion with a GP or respiratory clinic promptly. After thoracoscopy, follow written discharge advice and attend drain management reviews.

\textbf{Contact your local emergency services for:}
\begin{itemize}
    \item Sudden severe breathlessness, blue lips, or collapse after chest procedures
    \item Heavy bleeding from wounds or drains, or coughing large amounts of blood
    \item High fever with confusion or severe chest pain suggesting complicated infection
\end{itemize}

\section*{Diagnosis and Investigations}
Before thoracoscopy, teams complete staging of the clinical problem and fitness assessment.

\begin{itemize}
    \item \textbf{Imaging:} chest X-ray, ultrasound for effusion mapping, and CT chest; PET-CT when cancer staging is relevant
    \item \textbf{Thoracentesis:} protein, LDH, cytology, microbiology, and sometimes adenosine deaminase
    \item \textbf{Blood tests:} coagulation, blood count, renal function, and cross-match if bleeding risk
    \item \textbf{Lung function and anaesthetic review:} especially for VATS under general anaesthesia
    \item \textbf{Medication plan:} manage anticoagulants and antiplatelets with the proceduralist
    \item \textbf{Consent:} discussion of diagnostic yield, failure risk, pain, infection, bleeding, air leak, and rare need to convert to open thoracotomy
\end{itemize}

\section*{Management}
\begin{itemize}
    \item \textbf{Medical thoracoscopy:} lateral position, local anaesthetic, one or two ports, fluid removal, inspection, biopsy, and often talc poudrage; an intercostal catheter remains temporarily
    \item \textbf{VATS:} general anaesthesia, camera and instrument ports, definitive surgical tasks, stapling, and drain placement
    \item \textbf{Analgesia:} multimodal pain control including local anaesthetic techniques; good pain relief enables deep breathing and reduces pneumonia risk
    \item \textbf{Chest drain care:} underwater seal or digital drainage; removal when air leak stops and output is acceptable
    \item \textbf{Physiotherapy:} early mobilisation and breathing exercises
    \item \textbf{Histopathology follow-up:} results guide oncology, antimicrobial, or anti-inflammatory treatment
    \item \textbf{Indwelling pleural catheters:} alternative or adjunct for some malignant effusions when pleurodesis is unsuitable
\end{itemize}

\section*{Complications}
\begin{itemize}
    \item \textbf{Pain and intercostal neuralgia:} usually improves; occasionally persists
    \item \textbf{Bleeding:} may require return to theatre or transfusion
    \item \textbf{Infection and empyema:} uncommon with sterile technique but serious
    \item \textbf{Prolonged air leak and subcutaneous emphysema:} more likely with diseased lung
    \item \textbf{Re-expansion pulmonary oedema:} rare after rapid drainage of large chronic effusions
    \item \textbf{Anaesthetic and cardiac events:} related to comorbidity
    \item \textbf{Failed pleurodesis or non-diagnostic biopsy:} may need further procedures
    \item \textbf{Tumour seeding of port sites:} discussed particularly in mesothelioma pathways
\end{itemize}

\section*{Prognosis and Outcomes}
Diagnostic medical thoracoscopy yields a diagnosis in a high proportion of unexplained exudative effusions when performed by experienced operators. Talc pleurodesis succeeds in many malignant effusions, improving breathlessness control. VATS for pneumothorax greatly reduces recurrence compared with conservative care alone in appropriately selected patients. Lung cancer resection outcomes depend on stage and fitness. Recovery after uncomplicated thoracoscopy is often measured in days to a few weeks, with gradual return to full activity guided by the surgeon or physician.

\section*{Prevention and Self-Care}
\begin{itemize}
    \item Stop smoking before elective chest surgery to improve healing and reduce air leak and infection risk
    \item Follow fasting, medication, and transport instructions exactly on the procedure day
    \item Use incentive spirometry or prescribed breathing exercises after drains are in place
    \item Keep wound sites clean and dry as advised; report redness, pus, or fever early
    \item Do not remove or clamp drains yourself; seek the treating team if a drain dislodges
    \item Avoid heavy lifting and air travel until cleared, especially after pneumothorax treatment
    \item Attend all oncology or respiratory follow-up so biopsy results translate into treatment
\end{itemize}

\section*{Sources and Further Reading}
\begin{itemize}
    \item NHS. \textit{Thoracoscopy and VATS.} https://www.nhs.uk/
    \item Mayo Clinic. \textit{Thoracoscopy.} https://www.mayoclinic.org/
    \item MSD Manual. \textit{Thoracoscopy.} https://www.msdmanuals.com/
    \item MedlinePlus. \textit{Thoracoscopy.} https://medlineplus.gov/
    \item Better Health Channel. \textit{Chest procedures.} https://www.betterhealth.vic.gov.au/
    \item Healthdirect. \textit{Thoracoscopy.} https://www.healthdirect.gov.au/
\end{itemize}
